
\documentclass{article}
\usepackage{colortbl}
\usepackage{makecell}
\usepackage{multirow}
\usepackage{supertabular}

\begin{document}

\newcounter{utterance}

\centering \large Interaction Transcript for game `dond', experiment `semi\_en', episode 11 with Qwen3{-}32B{-}t0.0.
\vspace{24pt}

{ \footnotesize  \setcounter{utterance}{1}
\setlength{\tabcolsep}{0pt}
\begin{supertabular}{c@{$\;$}|p{.15\linewidth}@{}p{.15\linewidth}p{.15\linewidth}p{.15\linewidth}p{.15\linewidth}p{.15\linewidth}}
    \# & $\;$A & \multicolumn{4}{c}{Game Master} & $\;\:$B\\
    \hline

    \theutterance \stepcounter{utterance}  
    & & \multicolumn{4}{p{0.6\linewidth}}{
        \cellcolor[rgb]{0.9,0.9,0.9}{
            \makecell[{{p{\linewidth}}}]{
                \texttt{\tiny{[P1$\langle$GM]}}
                \texttt{You are playing a negotiation game in which you have to agree on how to divide a set of items among you and another player.} \\
\\ 
\texttt{Rules:} \\
\texttt{(a) You and the other player are given a set of items. Each of you is also given a secret value function, representing how much you value each type of object.} \\
\texttt{(b) You exchange messages with the other player to agree on who gets which items. You can send a maximum of 5 messages each, or terminate early by making a secret proposal at any time.} \\
\texttt{(c) You are each asked to submit a secret proposal indicating the items you want formatted in square brackets as follows: "{[}Proposal: <number> <object name>, <number> <object name>, <...>{]}"} \\
\texttt{(d) If your proposals are complementary, i.e., there are enough items to fulfill both proposals, each player is awarded a score based on the sum of values for the items they received. Otherwise, both of you get zero points.} \\
\texttt{(e) Your goal is to maximize the score you receive.} \\
\\ 
\texttt{Let us start.} \\
\\ 
\texttt{The set of available items is:} \\
\\ 
\texttt{1 pot, 1 candy, 1 whistle, 2 dice} \\
\\ 
\texttt{Your secret value function is:} \\
\\ 
\texttt{pot: 2, candy: 3, whistle: 3, dice: 1} \\
\\ 
\texttt{IMPORTANT: Your messages, unless it is the secret proposal, are directly transmitted to the other player, so do not include any response to the rules or text announcing your message. To make a secret proposal, use the indicated format. Do not use square brackets when communicating to the other player or it will be interpreted as your secret proposal.} \\
            }
        }
    }
    & \\ \\

    \theutterance \stepcounter{utterance}  
    & \multicolumn{4}{p{0.6\linewidth}}{
        \cellcolor[rgb]{0.9,0.9,0.9}{
            \makecell[{{p{\linewidth}}}]{
                \texttt{\tiny{[P1$\rangle$GM]}}
                \texttt{<think>} \\
\texttt{Okay, let's see. I need to negotiate with another player to divide the items: 1 pot, 1 candy, 1 whistle, and 2 dice. My value function is pot:2, candy:3, whistle:3, dice:1 each. So, my highest values are candy and whistle at 3 each, then pot at 2, and dice at 1 each. The other player's value function is secret, so I don't know what they value. } \\
\\ 
\texttt{First, I should figure out what the other player might want. Since the items are limited, maybe they also value some of the higher{-}value items. But without knowing their values, I have to make assumptions. Maybe they also want candy or whistle? Or maybe they prefer dice? Since there are two dice, maybe they can take both if they value them highly.} \\
\\ 
\texttt{My goal is to maximize my own score. So I need to get the items that give me the highest total. Let's think about possible splits. If I can get both candy and whistle, that's 3+3=6. Then the other player could take the pot and dice. But if they also want candy or whistle, that might not work. Alternatively, if I take one of the high{-}value items and some others.} \\
\\ 
\texttt{But since the proposals have to be complementary, I need to make sure that my proposal doesn't conflict with theirs. If I propose taking candy and whistle, and they also want those, then our proposals would conflict, and we both get zero. So maybe I should propose taking one high{-}value item and some others, leaving the other high{-}value item for them.} \\
\\ 
\texttt{Wait, but how do I know what they want? Since I can send up to 5 messages, maybe I can ask them what they value. But the problem is that they might lie or not respond. Alternatively, I can make an offer that's mutually beneficial. For example, if I take candy and one die, and they take whistle and pot and one die. Let's check the values. For me: candy (3) + die (1) = 4. For them: whistle (unknown value) + pot (unknown) + die (unknown). But since I don't know their values, maybe this isn't the best approach.} \\
\\ 
\texttt{Alternatively, maybe I should propose taking candy and whistle, which gives me 6 points. Then the remaining items are pot and two dice. If the other} \\
            }
        }
    }
    & & \\ \\

    \theutterance \stepcounter{utterance}  
    & & & \multicolumn{2}{p{0.3\linewidth}}{
        \cellcolor[rgb]{0.9,0.9,0.9}{
            \makecell[{{p{\linewidth}}}]{
                \texttt{\tiny{[GM$|$GM]}}
                \texttt{continue} \\
            }
        }
    }
    & & \\ \\

    \theutterance \stepcounter{utterance}  
    & & & \multicolumn{2}{p{0.3\linewidth}}{
        \cellcolor[rgb]{0.9,0.9,0.9}{
            \makecell[{{p{\linewidth}}}]{
                \texttt{\tiny{[GM$|$GM]}}
                \texttt{<think>} \\
\texttt{Okay, let's see. I need to negotiate with another player to divide the items: 1 pot, 1 candy, 1 whistle, and 2 dice. My value function is pot:2, candy:3, whistle:3, dice:1 each. So, my highest values are candy and whistle at 3 each, then pot at 2, and dice at 1 each. The other player's value function is secret, so I don't know what they value. } \\
\\ 
\texttt{First, I should figure out what the other player might want. Since the items are limited, maybe they also value some of the higher{-}value items. But without knowing their values, I have to make assumptions. Maybe they also want candy or whistle? Or maybe they prefer dice? Since there are two dice, maybe they can take both if they value them highly.} \\
\\ 
\texttt{My goal is to maximize my own score. So I need to get the items that give me the highest total. Let's think about possible splits. If I can get both candy and whistle, that's 3+3=6. Then the other player could take the pot and dice. But if they also want candy or whistle, that might not work. Alternatively, if I take one of the high{-}value items and some others.} \\
\\ 
\texttt{But since the proposals have to be complementary, I need to make sure that my proposal doesn't conflict with theirs. If I propose taking candy and whistle, and they also want those, then our proposals would conflict, and we both get zero. So maybe I should propose taking one high{-}value item and some others, leaving the other high{-}value item for them.} \\
\\ 
\texttt{Wait, but how do I know what they want? Since I can send up to 5 messages, maybe I can ask them what they value. But the problem is that they might lie or not respond. Alternatively, I can make an offer that's mutually beneficial. For example, if I take candy and one die, and they take whistle and pot and one die. Let's check the values. For me: candy (3) + die (1) = 4. For them: whistle (unknown value) + pot (unknown) + die (unknown). But since I don't know their values, maybe this isn't the best approach.} \\
\\ 
\texttt{Alternatively, maybe I should propose taking candy and whistle, which gives me 6 points. Then the remaining items are pot and two dice. If the other} \\
            }
        }
    }
    & & \\ \\

    \theutterance \stepcounter{utterance}  
    & & \multicolumn{4}{p{0.6\linewidth}}{
        \cellcolor[rgb]{0.9,0.9,0.9}{
            \makecell[{{p{\linewidth}}}]{
                \texttt{\tiny{[GM$\rangle$P2]}}
                \texttt{You are playing a negotiation game in which you have to agree on how to divide a set of items among you and another player.} \\
\\ 
\texttt{Rules:} \\
\texttt{(a) You and the other player are given a set of items. Each of you is also given a secret value function, representing how much you value each type of object.} \\
\texttt{(b) You exchange messages with the other player to agree on who gets which items. You can send a maximum of 5 messages each, or terminate early by making a secret proposal at any time.} \\
\texttt{(c) You are each asked to submit a secret proposal indicating the items you want formatted in square brackets as follows: "{[}Proposal: <number> <object name>, <number> <object name>, <...>{]}"} \\
\texttt{(d) If your proposals are complementary, i.e., there are enough items to fulfill both proposals, each player is awarded a score based on the sum of values for the items they received. Otherwise, both of you get zero points.} \\
\texttt{(e) Your goal is to maximize the score you receive.} \\
\\ 
\texttt{Let us start.} \\
\\ 
\texttt{The set of available items is:} \\
\\ 
\texttt{1 pot, 1 candy, 1 whistle, 2 dice} \\
\\ 
\texttt{Your secret value function is:} \\
\\ 
\texttt{pot: 8, candy: 0, whistle: 0, dice: 1} \\
\\ 
\texttt{IMPORTANT: Your messages, unless it is the secret proposal, are directly transmitted to the other player, so do not include any response to the rules or text announcing your message. To make a secret proposal, use the indicated format. Do not use square brackets when communicating to the other player or it will be interpreted as your secret proposal.} \\
\\ 
\\ 
\texttt{<think>} \\
\texttt{Okay, let's see. I need to negotiate with another player to divide the items: 1 pot, 1 candy, 1 whistle, and 2 dice. My value function is pot:2, candy:3, whistle:3, dice:1 each. So, my highest values are candy and whistle at 3 each, then pot at 2, and dice at 1 each. The other player's value function is secret, so I don't know what they value. } \\
\\ 
\texttt{First, I should figure out what the other player might want. Since the items are limited, maybe they also value some of the higher{-}value items. But without knowing their values, I have to make assumptions. Maybe they also want candy or whistle? Or maybe they prefer dice? Since there are two dice, maybe they can take both if they value them highly.} \\
\\ 
\texttt{My goal is to maximize my own score. So I need to get the items that give me the highest total. Let's think about possible splits. If I can get both candy and whistle, that's 3+3=6. Then the other player could take the pot and dice. But if they also want candy or whistle, that might not work. Alternatively, if I take one of the high{-}value items and some others.} \\
\\ 
\texttt{But since the proposals have to be complementary, I need to make sure that my proposal doesn't conflict with theirs. If I propose taking candy and whistle, and they also want those, then our proposals would conflict, and we both get zero. So maybe I should propose taking one high{-}value item and some others, leaving the other high{-}value item for them.} \\
\\ 
\texttt{Wait, but how do I know what they want? Since I can send up to 5 messages, maybe I can ask them what they value. But the problem is that they might lie or not respond. Alternatively, I can make an offer that's mutually beneficial. For example, if I take candy and one die, and they take whistle and pot and one die. Let's check the values. For me: candy (3) + die (1) = 4. For them: whistle (unknown value) + pot (unknown) + die (unknown). But since I don't know their values, maybe this isn't the best approach.} \\
\\ 
\texttt{Alternatively, maybe I should propose taking candy and whistle, which gives me 6 points. Then the remaining items are pot and two dice. If the other} \\
            }
        }
    }
    & \\ \\

    \theutterance \stepcounter{utterance}  
    & & & \multicolumn{4}{p{0.6\linewidth}}{
        \cellcolor[rgb]{0.9,0.9,0.9}{
            \makecell[{{p{\linewidth}}}]{
                \texttt{\tiny{[GM$\langle$P2]}}
                \texttt{<think>} \\
\texttt{Okay, let's see. The user is playing a negotiation game where they need to divide items between themselves and another player. The items are 1 pot, 1 candy, 1 whistle, and 2 dice. The user's value function is pot:2, candy:3, whistle:3, dice:1 each. The other player's values are secret. The goal is to maximize their own score by making a secret proposal that's complementary to the other player's.} \\
\\ 
\texttt{First, the user needs to figure out what the other player might value. Since the other player's values are unknown, the user has to make educated guesses. The user's highest values are candy and whistle (3 each), then pot (2), and dice (1 each). The other player might also value some of these items, but it's uncertain.} \\
\\ 
\texttt{The key here is to make a proposal that allows both players to get items without overlapping. If the user takes candy and whistle (total 6 points), the remaining items are pot and two dice. If the other player values pot highly, they might take that and one die, leaving the other die for the user. But the user only gets 1 die, which is 1 point, so their total would be 6 + 1 = 7? Wait, no. Wait, the user's proposal is to take candy and whistle, so they can't take the die unless it's part of their proposal. Wait, the proposal is a secret, so if the user proposes taking candy and whistle, and the other player proposes taking pot and two dice, then their proposals are complementary. The user's score would be 3 (candy) + 3 (whistle) = 6. The other player's score would be their value for pot plus two dice. But since the user doesn't know the other's values, they can't be sure. However, if the other player's values are such that they can take pot and two dice without conflicting, then the user's proposal would work.} \\
\\ 
\texttt{Alternatively, if the user takes candy, whistle, and one die, that's 3 + 3 + 1 = 7. Then the remaining items are pot and one die. The other player could take those. But if the other player's proposal includes taking candy or whistle, there's a conflict. So the user needs to make sure their proposal doesn't overlap with the other's. Since the other} \\
            }
        }
    }
     \\ \\

    \theutterance \stepcounter{utterance}  
    & & & \multicolumn{2}{p{0.3\linewidth}}{
        \cellcolor[rgb]{0.9,0.9,0.9}{
            \makecell[{{p{\linewidth}}}]{
                \texttt{\tiny{[GM$|$GM]}}
                \texttt{continue} \\
            }
        }
    }
    & & \\ \\

    \theutterance \stepcounter{utterance}  
    & & & \multicolumn{2}{p{0.3\linewidth}}{
        \cellcolor[rgb]{0.9,0.9,0.9}{
            \makecell[{{p{\linewidth}}}]{
                \texttt{\tiny{[GM$|$GM]}}
                \texttt{<think>} \\
\texttt{Okay, let's see. The user is playing a negotiation game where they need to divide items between themselves and another player. The items are 1 pot, 1 candy, 1 whistle, and 2 dice. The user's value function is pot:2, candy:3, whistle:3, dice:1 each. The other player's values are secret. The goal is to maximize their own score by making a secret proposal that's complementary to the other player's.} \\
\\ 
\texttt{First, the user needs to figure out what the other player might value. Since the other player's values are unknown, the user has to make educated guesses. The user's highest values are candy and whistle (3 each), then pot (2), and dice (1 each). The other player might also value some of these items, but it's uncertain.} \\
\\ 
\texttt{The key here is to make a proposal that allows both players to get items without overlapping. If the user takes candy and whistle (total 6 points), the remaining items are pot and two dice. If the other player values pot highly, they might take that and one die, leaving the other die for the user. But the user only gets 1 die, which is 1 point, so their total would be 6 + 1 = 7? Wait, no. Wait, the user's proposal is to take candy and whistle, so they can't take the die unless it's part of their proposal. Wait, the proposal is a secret, so if the user proposes taking candy and whistle, and the other player proposes taking pot and two dice, then their proposals are complementary. The user's score would be 3 (candy) + 3 (whistle) = 6. The other player's score would be their value for pot plus two dice. But since the user doesn't know the other's values, they can't be sure. However, if the other player's values are such that they can take pot and two dice without conflicting, then the user's proposal would work.} \\
\\ 
\texttt{Alternatively, if the user takes candy, whistle, and one die, that's 3 + 3 + 1 = 7. Then the remaining items are pot and one die. The other player could take those. But if the other player's proposal includes taking candy or whistle, there's a conflict. So the user needs to make sure their proposal doesn't overlap with the other's. Since the other} \\
            }
        }
    }
    & & \\ \\

    \theutterance \stepcounter{utterance}  
    & & \multicolumn{4}{p{0.6\linewidth}}{
        \cellcolor[rgb]{0.9,0.9,0.9}{
            \makecell[{{p{\linewidth}}}]{
                \texttt{\tiny{[P1$\langle$GM]}}
                \texttt{<think>} \\
\texttt{Okay, let's see. The user is playing a negotiation game where they need to divide items between themselves and another player. The items are 1 pot, 1 candy, 1 whistle, and 2 dice. The user's value function is pot:2, candy:3, whistle:3, dice:1 each. The other player's values are secret. The goal is to maximize their own score by making a secret proposal that's complementary to the other player's.} \\
\\ 
\texttt{First, the user needs to figure out what the other player might value. Since the other player's values are unknown, the user has to make educated guesses. The user's highest values are candy and whistle (3 each), then pot (2), and dice (1 each). The other player might also value some of these items, but it's uncertain.} \\
\\ 
\texttt{The key here is to make a proposal that allows both players to get items without overlapping. If the user takes candy and whistle (total 6 points), the remaining items are pot and two dice. If the other player values pot highly, they might take that and one die, leaving the other die for the user. But the user only gets 1 die, which is 1 point, so their total would be 6 + 1 = 7? Wait, no. Wait, the user's proposal is to take candy and whistle, so they can't take the die unless it's part of their proposal. Wait, the proposal is a secret, so if the user proposes taking candy and whistle, and the other player proposes taking pot and two dice, then their proposals are complementary. The user's score would be 3 (candy) + 3 (whistle) = 6. The other player's score would be their value for pot plus two dice. But since the user doesn't know the other's values, they can't be sure. However, if the other player's values are such that they can take pot and two dice without conflicting, then the user's proposal would work.} \\
\\ 
\texttt{Alternatively, if the user takes candy, whistle, and one die, that's 3 + 3 + 1 = 7. Then the remaining items are pot and one die. The other player could take those. But if the other player's proposal includes taking candy or whistle, there's a conflict. So the user needs to make sure their proposal doesn't overlap with the other's. Since the other} \\
            }
        }
    }
    & \\ \\

    \theutterance \stepcounter{utterance}  
    & \multicolumn{4}{p{0.6\linewidth}}{
        \cellcolor[rgb]{0.9,0.9,0.9}{
            \makecell[{{p{\linewidth}}}]{
                \texttt{\tiny{[P1$\rangle$GM]}}
                \texttt{<think>} \\
\texttt{Okay, let's try to figure out the best way to approach this. The user has to divide the items between themselves and another player, with their own value function. The items are 1 pot, 1 candy, 1 whistle, and 2 dice. The user's values are pot:2, candy:3, whistle:3, dice:1 each. The other player's values are secret, so the user has to make a proposal that's likely to be complementary.} \\
\\ 
\texttt{First, the user's highest values are candy and whistle at 3 each. So ideally, they want those. But if they take both, the other player would have to take the pot and the two dice. But the other player might also value candy or whistle, so maybe they would propose taking one of those. However, since the user doesn't know the other's values, they have to make an assumption. If the user takes candy and whistle, their total would be 6. The remaining items are pot and two dice. If the other player can take those without overlapping, then it's a valid proposal. But if the other player also wants candy or whistle, their proposals would conflict.} \\
\\ 
\texttt{Alternatively, maybe the user should take one of the high{-}value items and some others. For example, take candy and one die (3 + 1 = 4) and leave the whistle for the other player. Then the other player could take whistle, pot, and the other die. That way, the user gets 4, and the other player gets their values for whistle, pot, and die. But the user's total would be lower than if they took both candy and whistle. However, this might be safer if the other player also wants one of the high{-}value items.} \\
\\ 
\texttt{Another option is to take the pot and both dice. The user's value would be 2 + 1 + 1 = 4. Then the other player could take candy and whistle, which would give them 3 + 3 = 6. But the user's own score would be lower. But maybe the other player would prefer that, so the user's proposal would be complementary. However, the user's goal is to maximize their own score, so this might not be ideal.} \\
\\ 
\texttt{Wait, the user's maximum possible score is if they get both candy and whistle (6 points). The other items (pot and dice) are lower value. So the best strategy is to try to} \\
            }
        }
    }
    & & \\ \\

    \theutterance \stepcounter{utterance}  
    & & & \multicolumn{2}{p{0.3\linewidth}}{
        \cellcolor[rgb]{0.9,0.9,0.9}{
            \makecell[{{p{\linewidth}}}]{
                \texttt{\tiny{[GM$|$GM]}}
                \texttt{continue} \\
            }
        }
    }
    & & \\ \\

    \theutterance \stepcounter{utterance}  
    & & & \multicolumn{2}{p{0.3\linewidth}}{
        \cellcolor[rgb]{0.9,0.9,0.9}{
            \makecell[{{p{\linewidth}}}]{
                \texttt{\tiny{[GM$|$GM]}}
                \texttt{<think>} \\
\texttt{Okay, let's try to figure out the best way to approach this. The user has to divide the items between themselves and another player, with their own value function. The items are 1 pot, 1 candy, 1 whistle, and 2 dice. The user's values are pot:2, candy:3, whistle:3, dice:1 each. The other player's values are secret, so the user has to make a proposal that's likely to be complementary.} \\
\\ 
\texttt{First, the user's highest values are candy and whistle at 3 each. So ideally, they want those. But if they take both, the other player would have to take the pot and the two dice. But the other player might also value candy or whistle, so maybe they would propose taking one of those. However, since the user doesn't know the other's values, they have to make an assumption. If the user takes candy and whistle, their total would be 6. The remaining items are pot and two dice. If the other player can take those without overlapping, then it's a valid proposal. But if the other player also wants candy or whistle, their proposals would conflict.} \\
\\ 
\texttt{Alternatively, maybe the user should take one of the high{-}value items and some others. For example, take candy and one die (3 + 1 = 4) and leave the whistle for the other player. Then the other player could take whistle, pot, and the other die. That way, the user gets 4, and the other player gets their values for whistle, pot, and die. But the user's total would be lower than if they took both candy and whistle. However, this might be safer if the other player also wants one of the high{-}value items.} \\
\\ 
\texttt{Another option is to take the pot and both dice. The user's value would be 2 + 1 + 1 = 4. Then the other player could take candy and whistle, which would give them 3 + 3 = 6. But the user's own score would be lower. But maybe the other player would prefer that, so the user's proposal would be complementary. However, the user's goal is to maximize their own score, so this might not be ideal.} \\
\\ 
\texttt{Wait, the user's maximum possible score is if they get both candy and whistle (6 points). The other items (pot and dice) are lower value. So the best strategy is to try to} \\
            }
        }
    }
    & & \\ \\

    \theutterance \stepcounter{utterance}  
    & & \multicolumn{4}{p{0.6\linewidth}}{
        \cellcolor[rgb]{0.9,0.9,0.9}{
            \makecell[{{p{\linewidth}}}]{
                \texttt{\tiny{[GM$\rangle$P2]}}
                \texttt{<think>} \\
\texttt{Okay, let's try to figure out the best way to approach this. The user has to divide the items between themselves and another player, with their own value function. The items are 1 pot, 1 candy, 1 whistle, and 2 dice. The user's values are pot:2, candy:3, whistle:3, dice:1 each. The other player's values are secret, so the user has to make a proposal that's likely to be complementary.} \\
\\ 
\texttt{First, the user's highest values are candy and whistle at 3 each. So ideally, they want those. But if they take both, the other player would have to take the pot and the two dice. But the other player might also value candy or whistle, so maybe they would propose taking one of those. However, since the user doesn't know the other's values, they have to make an assumption. If the user takes candy and whistle, their total would be 6. The remaining items are pot and two dice. If the other player can take those without overlapping, then it's a valid proposal. But if the other player also wants candy or whistle, their proposals would conflict.} \\
\\ 
\texttt{Alternatively, maybe the user should take one of the high{-}value items and some others. For example, take candy and one die (3 + 1 = 4) and leave the whistle for the other player. Then the other player could take whistle, pot, and the other die. That way, the user gets 4, and the other player gets their values for whistle, pot, and die. But the user's total would be lower than if they took both candy and whistle. However, this might be safer if the other player also wants one of the high{-}value items.} \\
\\ 
\texttt{Another option is to take the pot and both dice. The user's value would be 2 + 1 + 1 = 4. Then the other player could take candy and whistle, which would give them 3 + 3 = 6. But the user's own score would be lower. But maybe the other player would prefer that, so the user's proposal would be complementary. However, the user's goal is to maximize their own score, so this might not be ideal.} \\
\\ 
\texttt{Wait, the user's maximum possible score is if they get both candy and whistle (6 points). The other items (pot and dice) are lower value. So the best strategy is to try to} \\
            }
        }
    }
    & \\ \\

    \theutterance \stepcounter{utterance}  
    & & & \multicolumn{4}{p{0.6\linewidth}}{
        \cellcolor[rgb]{0.9,0.9,0.9}{
            \makecell[{{p{\linewidth}}}]{
                \texttt{\tiny{[GM$\langle$P2]}}
                \texttt{<think>} \\
\texttt{Okay, let's see. The user is trying to maximize their score by making a secret proposal. The items are 1 pot, 1 candy, 1 whistle, and 2 dice. Their value function is pot:2, candy:3, whistle:3, dice:1 each. The other player's values are unknown.} \\
\\ 
\texttt{The user's highest values are candy and whistle at 3 each. So ideally, they want those. If they take both, their total would be 6. The remaining items are pot and two dice. The other player would need to take those. But if the other player also wants candy or whistle, their proposals would conflict. However, since the user doesn't know the other's values, they have to assume that the other player might not want those items. Alternatively, maybe the other player values the pot or dice more.} \\
\\ 
\texttt{But the user's goal is to maximize their own score. So taking candy and whistle gives them 6 points. The other player would get pot and two dice. If the other player's values for those items are such that they can take them without overlapping, then the proposals are complementary. Since the user can't know for sure, but the rules state that if the proposals are complementary, they get their scores. Otherwise, both get zero. So the user needs to make a proposal that is likely to be complementary.} \\
\\ 
\texttt{Another option is to take one high{-}value item and some others. For example, take candy and one die (3 + 1 = 4) and leave the whistle for the other player. Then the other player could take whistle, pot, and the other die. That way, the user's score is 4, but the other player might get more. But the user's own score is lower. However, this might be safer if the other player also wants one of the high{-}value items.} \\
\\ 
\texttt{Alternatively, the user could take the pot and both dice (2 + 1 + 1 = 4) and leave candy and whistle for the other player. But that would give the user only 4 points, which is worse than taking candy and whistle.} \\
\\ 
\texttt{So the best strategy seems to be to take candy and whistle (6 points) and let the other player take pot and two dice. The user's proposal would be {[}Proposal: 1 candy, 1 whistle{]}. The other player's proposal would need to be for the remaining items. If the other player's proposal} \\
            }
        }
    }
     \\ \\

    \theutterance \stepcounter{utterance}  
    & & & \multicolumn{2}{p{0.3\linewidth}}{
        \cellcolor[rgb]{0.9,0.9,0.9}{
            \makecell[{{p{\linewidth}}}]{
                \texttt{\tiny{[GM$|$GM]}}
                \texttt{proposal} \\
            }
        }
    }
    & & \\ \\

    \theutterance \stepcounter{utterance}  
    & & & \multicolumn{2}{p{0.3\linewidth}}{
        \cellcolor[rgb]{0.9,0.9,0.9}{
            \makecell[{{p{\linewidth}}}]{
                \texttt{\tiny{[GM$|$GM]}}
                [0, 1, 1, 0]
            }
        }
    }
    & & \\ \\

    \theutterance \stepcounter{utterance}  
    & & \multicolumn{4}{p{0.6\linewidth}}{
        \cellcolor[rgb]{0.9,0.9,0.9}{
            \makecell[{{p{\linewidth}}}]{
                \texttt{\tiny{[P1$\langle$GM]}}
                \texttt{The other player has submitted a secret proposal. It is now your turn to submit a proposal of your own.} \\
            }
        }
    }
    & \\ \\

    \theutterance \stepcounter{utterance}  
    & \multicolumn{4}{p{0.6\linewidth}}{
        \cellcolor[rgb]{0.9,0.9,0.9}{
            \makecell[{{p{\linewidth}}}]{
                \texttt{\tiny{[P1$\rangle$GM]}}
                \texttt{<think>} \\
\texttt{Okay, the other player has already submitted a secret proposal, and now it's my turn. I need to make sure my proposal is complementary to theirs. Since I don't know their proposal, I have to make an educated guess based on the items and my own value function.} \\
\\ 
\texttt{My value function is pot:2, candy:3, whistle:3, dice:1 each. The available items are 1 pot, 1 candy, 1 whistle, and 2 dice. The other player's values are unknown, but I need to maximize my own score.} \\
\\ 
\texttt{If I take the candy and whistle, that gives me 3 + 3 = 6 points. The remaining items would be pot and two dice. If the other player's proposal includes taking pot and the two dice, then our proposals are complementary. But if they also wanted candy or whistle, there's a conflict. However, since I don't know their values, I have to assume they might not need those. Alternatively, maybe they want one of the high{-}value items, so I should leave one for them.} \\
\\ 
\texttt{Another option is to take one high{-}value item and some others. For example, take candy and one die (3 + 1 = 4) and leave the whistle for the other player. Then the other player could take whistle, pot, and the other die. That way, my score is 4, and theirs would be their values for those items. But my score is lower than if I took both candy and whistle.} \\
\\ 
\texttt{Alternatively, if I take the pot and both dice (2 + 1 + 1 = 4), then the other player could take candy and whistle (6 points). But again, my score is lower. However, if the other player's proposal is for candy and whistle, then mine would conflict. Wait, but if I take pot and dice, and they take candy and whistle, that's complementary. So maybe that's a safer bet if I think they might want the high{-}value items.} \\
\\ 
\texttt{But since my goal is to maximize my own score, I should aim for the highest possible. Taking candy and whistle gives me 6, which is the maximum possible. The other items are pot and two dice. If the other player's proposal is for pot and dice, then it's okay. If they also wanted candy or whistle, then we both get zero. But without knowing their values, I have to take the risk. Since the other} \\
            }
        }
    }
    & & \\ \\

    \theutterance \stepcounter{utterance}  
    & & & \multicolumn{2}{p{0.3\linewidth}}{
        \cellcolor[rgb]{0.9,0.9,0.9}{
            \makecell[{{p{\linewidth}}}]{
                \texttt{\tiny{[GM$|$GM]}}
                \texttt{continue} \\
            }
        }
    }
    & & \\ \\

    \theutterance \stepcounter{utterance}  
    & & & \multicolumn{2}{p{0.3\linewidth}}{
        \cellcolor[rgb]{0.9,0.9,0.9}{
            \makecell[{{p{\linewidth}}}]{
                \texttt{\tiny{[GM$|$GM]}}
                \texttt{player was instructed to make a proposal} \\
            }
        }
    }
    & & \\ \\

\end{supertabular}
}

\end{document}
